%% =====================================================================
%%  T-17-01  Master Timing
%%  -------------------------------------------------------------------
%%  One idea per file. Core is mandatory. Slots in canonical order.
%%  Max 700 words. No layout commands. No filler. New source material
%%  for this entry goes in sources/<BOOK>/T-17-01.tex -- never by
%%  rewriting this file (docs/RULES.md R0.1).
%%  =====================================================================
%% @kind: topic
%% @id: T-17-01
%% @title: Master Timing
%% @subtitle: The moment is part of the move
%% @chapter: c17
%% @order: 10
%% @status: stable
%% @sources: B3
%% @updated: 2026-09-19
%% @allow-rewrite: B3 ingest 2026-09-19 -- committed scaffold replaced by the written entry (planned -> stable lifecycle)

\topic{T-17-01}{Master Timing}{The moment is part of the move}

\begin{core}
The same move lands or dies by when it is made. Hurry betrays a lack of control and re-prices everything you ask for; patience reads as command, because you seem to know that what you want is coming to you. The craft is to sniff the spirit of the times, stand back while the moment is unripe, and strike fiercely when it ripens --- then stop.
\end{core}
\begin{mechanism}
Two clocks run at once. The first is yours: visible hurry signals that you need the outcome more than the counterpart does, which is a pricing signal against you in every exchange (T-02-03 is this entry's patience-as-position twin). Appearing unhurried is itself leverage --- desire concealed is desire not taxed. The second clock is the world's: every idea, style and campaign has a season, and reading it --- \emph{becoming a detective of the right moment} --- tells you which moves the present will reward. The same request, plan or reform lands differently in different weather; the skilful wait for the confluence of ripeness and readiness, then act with full force, because a ripe moment struck softly is wasted and an unripe one struck hard is buried. The planning law pairs with it: end with the end in view and execute along that line, so that timing serves a design rather than an itch.
\end{mechanism}
\begin{conditions}
Strongest in institutions, markets and relationships where attention comes in seasons --- budget cycles, crises, moods, openings after failures. Weakens in continuous, low-season domains (routine operations) where waiting is just delay, and against counterparts who deliberately create false seasons to move you on their clock.
\end{conditions}
\begin{application}
  \item Before any ask, place it in a season: is now the world's moment for this, or only yours?
  \item Remove hurry signals --- deadlines you volunteer, reminders you chase, the second follow-up inside a day.
  \item When unripe, prepare visibly and press invisibly; let the counterpart conclude the idea was theirs when the season turns.
  \item At ripeness, commit fully and fast. Half-strokes at the right moment lose to full strokes.
  \item Declare the move finished when the objective is reached; the extra lap re-opens what closing settled (T-09-03).
\end{application}
\begin{feedback}
  \item Landing: the counterpart brings the subject up themselves --- the moment ripened on their side.
  \item Landing: your unhurried pace draws better terms without a word about terms.
  \item Failing: every nudge costs goodwill and moves nothing. The season is wrong; stop paying.
  \item Failing: you struck at ripeness and kept striking. The win is aging into an imposition.
\end{feedback}
\begin{failure}
Patience decays into procrastination wearing strategy as a coat --- waiting that has no ripeness condition is just fear with a calendar. Over-reading the times produces paralysis by analysis of moods. And the strike itself is exposed: the person who waits for perfect ripeness often finds the moment taken by someone who moved at good-enough.
\end{failure}
\begin{countermeasures}
  \item Put a date on your own patience. Open-ended waiting is someone else's clock wearing your clothes.
  \item Notice manufactured urgency from the other side --- ripeness announced by the party who benefits is a hook (T-24-03).
  \item In your own house, watch for the always-waiting plan. A strategy that never ripens is a comfort object.
  \item Give the unhurried counterpart a real season, not a fake deadline: genuine scarcity moves them, invented scarcity unmasks you.
\end{countermeasures}

\sources{B3}
\seesources
\seealso{T-02-03, T-17-02, T-24-03}
